\subsection{3.5 将多种能力整合进单一模型}\label{consolidating-capabilities-into-a-single-model}

前面的章节描述了三个教师模型，它们分别针对 STEM 与竞赛编程（第 3.2 节）、智能体能力（第 3.3 节）以及帮助性与安全性（第 3.4 节）独立训练。我们分两个阶段将它们整合进单一模型，如图 12 所示。SFT 阶段复用第 3.1.4 节的自蒸馏管线，将其应用于每个专项教师模型，不过三个教师模型需要不同的过滤和拒绝采样策略。对于 STEM 与智能体教师模型，我们按照第 3.1.4 节的多样性发现，在每次攀升的多个检查点上采样推演，并优先使用较晚的检查点。我们为每个上下文保留多条正确推演，只施加轻度过滤以剔除退化的思维链（CoT）。对于帮助性与安全性教师模型，我们使用 LLM 裁判和启发式过滤器，在正确性之外，还按风格、结构和已知缺陷对轨迹评分。

整合 SFT。SFT 阶段将三个教师模型蒸馏进单一模型。我们对数据混合配比和超参数进行了迭代，以平衡推理、智能体与通用帮助性能力，并在公开基准、内部基准和人工评估上进行验证。表 10 按样本权重和 token 权重报告了最终混合配比。我们发现按样本权重平衡混合配比很重要；相应地，token 分布由 STEM 和编程主导，这是因为它们的轨迹更长，而这在实践中并不会损害帮助性与安全性能力。与我们的标准自蒸馏方案不同，整合 SFT 训练 4 个 epoch，学习率从最大值 $1\cdot10^{-5}$ 衰减 2$\times$。

整合 RL。最后阶段的轻量 RL 进一步改进安全性、过度拒答和风格。该方案基于帮助性与安全性攀升（第 3.4 节），为保持推理性能做了一些调整。我们以 128k tokens 的最大序列长度训练，并在 RL 混合数据中保留少量 STEM 和编程数据；我们发现两者都很重要，否则复杂任务上的推理性能会在攀升过程中缓慢退化。

\subsection{3.6 强化学习基础设施}\label{reinforcement-learning-infrastructure}

我们的 RL 攀升依赖 Rocket，一个用于大规模异步分布式强化学习的内部框架，使用 YOLO 作为学习器（第 2.8 节），使用 SGLang（Zheng et al., 2024）进行模型推理。我们构建 Rocket，是因为我们的需求——即在 MAI-Thinking-1 所需的数千块 GPU 规模上支持异步 RL——是现有开源 RL 框架（von Werra et al., 2020; Hu et al., 2024; Sheng et al., 2024）无法满足的。Rocket 的核心 RL 数据流围绕单个控制器、一组问题工作节点和推演工作节点，以及生成模型输出的路由与推理服务器来组织。控制器、问题工作节点和推演工作节点各自实现为单个 Python 进程 / Ray actor，而路由与推理由基于 SGLang 的服务栈提供。Rocket 架构的总览如图 19 所示。



\pandocbounded{\includegraphics[keepaspectratio,alt={image}]{images/2b1dd5df63022b27e39ca9ddc41d8ad4d18b50e378b9c253e82dd35de310075e.jpg}}

图 19. Rocket 总览：我们的大规模强化学习框架。

\subsection{3.6.1 控制器}\label{controller}

控制器加载 RL 任务并将其发送给问题工作节点处理。对于每个任务，控制器会收到一组已完成的推演及关联的评分元数据（通过/失败判定、奖励、归一化优势等）。控制器对完成的推演进行过滤，并分批发送给学习器。控制器的抽象使 Rocket 可以在 on-policy（同策略）RL 与 off-policy（离策略）RL 之间切换不同的控制器实现。实践中，我们在大规模运行中主要使用 off-policy（离策略）RL，而将 on-policy（同策略）RL 保留给较小的实验和调试。

\subsection{3.6.2 问题工作节点}\label{problem-worker}

当问题工作节点从控制器收到 RL 任务时，它生成一组推演、为每条推演计算归一化优势，并将结果发回控制器。问题工作节点本身并不负责推演生成，而是向推演工作节点发送一组请求——每条推演一个请求。问题工作节点为大规模异步 RL 实现容错，例如当某个推演工作节点失败时为任务重试推演。

单个请求到达问题工作节点的典型流程如下：

\begin{enumerate}
\def\labelenumi{\arabic{enumi}.}
\tightlist
\item
  早退出推演。按照第 3.1.3 节描述的早退出策略，问题工作节点从控制器接收 RL 任务，并向推演工作节点发出 16 个请求，必要时重试失败的请求。每条推演都会被评分，这些信息被聚合成该任务的整体通过率。如果整体通过率落在预设区间内，问题工作节点进入完整推演阶段；否则，该 RL 任务被中止。
\end{enumerate}



\begin{enumerate}
\def\labelenumi{\arabic{enumi}.}
\setcounter{enumi}{1}
\item
  完整推演。问题工作节点再向推演工作节点发出 128 个请求，同样在必要时重试失败的请求。每条推演都会被评分，这些信息被聚合成该任务的整体通过率。如果整体通过率落在预设区间内，问题工作节点进入后处理阶段；否则，该 RL 任务被中止。
\item
  后处理。问题工作节点为每条推演计算归一化优势，用于 RL 目标（第 3.1.1 节）。视具体问题而定，该步骤可能包括对长推演施加长度惩罚、应用 GRPO 奖励归一化或其他启发式规则。
\end{enumerate}

\subsection{3.6.3 推演工作节点}\label{rollout-worker}

当推演工作节点收到来自问题工作节点的请求时，它生成一条推演，可选地进行评分，然后将结果发回调用方。在典型流程中，推演工作节点首先生成初始模型提示词并发送给模型推理服务器。响应返回后，推演工作节点解析响应、执行任何工具调用，并将结果合并进发给推理服务器的后续请求中。该过程持续进行，直到模型停止调用工具并输出最终答案，或超过预先设定的步数或推演耗时上限。

推演结束后，下一步是对模型的输出进行评分。评分器的输出包含数值奖励和可选的通过/失败判定。视问题而定，评分可以在推演工作节点上或问题工作节点上进行。如果各推演可以独立评估，评分由推演工作节点执行。例如，对于数学问题，如果模型的输出与问题的标准答案在数学上等价，评分器可能输出 (1, pass)，否则输出 (0, fail)。另一方面，有些评分器需要同时查看多条推演的答案，例如让 LLM 对不同模型的答案进行两两比较。在这种情况下，评分在问题工作节点上进行。单个问题可以有多个评分，此时由用户定义的策略将各评分聚合成一个整体奖励。

\subsection{3.6.4 路由与推理}\label{router-and-inference}

推理是系统最重要的组件之一：视任务的规模和特点而定，推理与学习器的 GPU 之比最高可达 5:1。在我们最大的 RL 任务中，共使用 4864 块 GB300 芯片，其中 4096 块专用于推理，其余 768 块专用于学习器。因此，推理性能和系统稳定性是一级关注点。为保障 RL 训练稳定性，控制学习器与推理引擎之间的数值差异至关重要。

推理配置。我们的 RL 推理栈构建在 SGLang 和 SGLang router（Zheng et al., 2024）之上。我们选择 SGLang，是因为它满足了我们对丰富 RL 服务特性、面向开源模型的快速实验以及强劲生态发展速度的需求。在单个 SGLang worker 之上，我们使用 SGLang router 负责负载均衡、流量控制和前缀缓存。

推理性能。由于任务的大部分 GPU 用量来自推理，优化推理吞吐量对高效的 RL 至关重要。最小化请求延迟对控制陈旧性同样关键。我们根据推演发生在单轮还是多轮场景来定制优化。对于单轮工作负载，提示词通常较短，但生成（包括思考）最长可达 128k tokens，因此主要瓶颈是 KV 缓存内存。我们通过专家并行和注意力数据并行来降低 KV 缓存与权重的占用。我们还禁用前缀缓存，使滑动窗口注意力 token 在长生成过程中可以被完全逐出（我们的模型架构重度依赖滑动窗口注意力）。我们还通过 MLP 层的数据并行、DeepEP（Deep-EveryParallel）（Zhao et al., 2025a）和 EPLB（Expert Parallelism Load Balancer）（DeepSeek-AI, 2025）进一步降低通信开销。对于多轮工作负载，提示词会变得很长，而生成往往较短，使这类负载以 prefill 为主。在这里我们重度依赖前缀缓存：在生产 RL 运行中，我们的前缀缓存命中率可达 97--98\%。



推理系统稳定性。在用于训练 MAI-Thinking-1 的、由数千块推理芯片构成的 RL 规模下，单个副本会时常因节点故障而崩溃、挂起、变慢或失败；此时推理层必须优雅降级并重启，而不是让整个任务失败。我们依赖三个层面的纵深防御。在副本层面，每个 SGLang server 运行一个自看门狗，探测自身的生成端点并监控调度器内存；不健康的副本会触发自身的优雅重启，使编排器可以拉起一个干净的新副本。在路由层面，SGLang router 充当熔断器：当后端副本不健康时，它停止接收新请求，重新开放前执行多阶段探活，并定期重新发现那些悄然退出其活跃集合的副本。逐副本流量控制可防止慢副本积压请求，从而拖累整个推理池。在任务层面，存活监控器跟踪每类 actor 的存活数量（推理副本、router、推演工作节点、学习器 rank），如果任何一类低于允许阈值，就让任务失败以进行干净的重启。另一个正交的步进进度看门狗捕获更棘手的情况：一切看起来都活着，但训练已停止推进。最后，当推演工作节点对推理服务器的请求在半途失败时，会针对另一个副本重试，因此单个副本故障只会降低推理池吞吐量，而不会丢失推演。

推理与学习器之间的数值差异。RL 的一个关键方面是推理引擎与学习器之间的数值差异。YOLO（学习器）和 SGLang（推理引擎）使用不同的 kernel、调度和并行策略。即使是微小的逐 token logprob 差异，也会在长推演中累积，并可能破坏 off-policy（离策略）RL 中的重要性采样校正。为缓解这些问题，我们在 RL 运行中对学习器和推理引擎都使用 bf16；与我们评估过的更低精度方案相比，这产生了显著更小的数值差异和更稳定的训练。在 bf16 之上，我们还应用 MoE 路由重放（Ma et al., 2025）和 top-p 掩码重放（第 3.1.3 节）。

\subsection{3.6.5 权重迁移}\label{weight-transfer}

从同步 RL 转向异步 RL，使权重迁移变成一个反复出现的分布式同步问题：新的学习器权重必须每隔 k 步到达推理集群，而迁移时间会造成推理空闲和额外的陈旧性。核心难点在于，学习器与推理的分配以不同的方式对同一批张量进行分片：双方可能使用不同组合或程度的 FSDP、流水线并行、数据并行注意力与张量布局，参数还可能在精度、量化状态或矩阵布局上不同。因此，每次迁移都必须既重新分片字节，又执行逐张量变换。

迁移计划。我们将重新分片、数据搬移和逐张量变换编译成单一的迁移计划，在任务开始时计算一次，之后复用。对每个参数，规划器求学习器与推理服务器分片布局的交集，并为每个非空的子分片重叠输出一条记录，记录源 rank、目标 rank、字节范围以及 dtype 转换、布局置换等所需操作。重新分片隐含在这份交集中；每个学习器 rank 只发送与其配对的推理 rank 所需的切片，从而省去对完整张量的单独物化。变换在成本最低的一侧执行，运行时将打包、传输和解包流水线化，使连续的多个子分片在时间上重叠。为了将规划与集群规模解耦，该计划面向一台学习器加一台推理服务器的理想拓扑；在运行时，每个源--目标对会扩展为一个覆盖目标 rank 全部存活副本的迁移组，因此当副本加入、离开或被替换时，计划仍然有效。由于数据并行是复制而非分片参数，一个 DP 组就已持有全部模型参数——计划无需征用全部 D 个学习器 rank。我们将其限制在部分 DP 组上，每个 DP 组与一段互不相交的推理副本配对并独立运行。一个 36 台服务器的集群按四路切分后，并行执行四次 9 台服务器的迁移，在提升吞吐量的同时控制故障爆炸半径。



表 11. 训练后模型在 STEM 与智能体编码公开基准上的评测结果。其他模型的数据取自各自的官方模型卡。得分以 ``---'' 列出的模型未报告相应基准的数值。对于智能体编码评测，我们使用 256k 的总上下文长度。对于上述其他评测，我们使用 256k 的最大输出 token 数。

{\def\LTcaptype{none} % do not increment counter
\begin{longtable}[]{@{}lllllllll@{}}
\toprule\noalign{}
\endhead
\bottomrule\noalign{}
\endlastfoot
\multirow{2}{*}{基准} & \multicolumn{8}{l@{}}{%
MAI-Thinking-1 Sonnet 4.6 Opus 4.6 GPT 5.4 Kimi K2.6 DeepSeek V3.2
DeepSeek V4 GLM-5.1} \\
AIME 2025 & 97.0 & 95.6 & 99.8 & --- & --- & 93.1 & --- & --- \\
AIME 2026 & 94.5 & --- & --- & --- & 96.4 & --- & --- & 95.3 \\
HMMT Feb 2026 & 84.9 & --- & --- & --- & 92.7 & --- & 95.2 & 82.6 \\
GPQA Diamond & 84.2 & 89.9 & 91.3 & 92.8 & 90.5 & 82.4 & 90.1 & 86.2 \\
LCB v6 & 87.7 & --- & --- & --- & 89.6 & 83.3 & 93.5 & --- \\
Terminal-Bench 2.0 & 46.0 & 59.1 & 65.4 & 75.1 & 66.7 & 46.4 & 67.9 &
69.0 \\
SWE-bench Verified & 73.5 & 79.6 & 80.8 & --- & 80.2 & 73.1 & 80.6 & --- \\
SWE-Bench Pro & 52.8 & --- & 53.4 & 57.7 & 58.6 & --- & 55.4 & 58.4 \\
\end{longtable}
}
